\IfFileExists{stacks-project.cls}{%
\documentclass{stacks-project}
}{%
\documentclass{amsart}
}

% For dealing with references we use the comment environment
\usepackage{verbatim}
\newenvironment{reference}{\comment}{\endcomment}
%\newenvironment{reference}{}{}
\newenvironment{slogan}{\comment}{\endcomment}
\newenvironment{history}{\comment}{\endcomment}

% For commutative diagrams we use Xy-pic
\usepackage[all]{xy}

% We use 2cell for 2-commutative diagrams.
\xyoption{2cell}
\UseAllTwocells

% We use multicol for the list of chapters between chapters
\usepackage{multicol}

% This is generally recommended for better output
\usepackage{lmodern}
\usepackage[T1]{fontenc}

% For cross-file-references
\usepackage{xr-hyper}

% Package for hypertext links:
\usepackage{hyperref}

% For any local file, say "hello.tex" you want to link to please
% use \externaldocument[hello-]{hello}
\externaldocument[introduction-]{introduction}
\externaldocument[conventions-]{conventions}
\externaldocument[sets-]{sets}
\externaldocument[categories-]{categories}
\externaldocument[topology-]{topology}
\externaldocument[sheaves-]{sheaves}
\externaldocument[sites-]{sites}
\externaldocument[stacks-]{stacks}
\externaldocument[fields-]{fields}
\externaldocument[algebra-]{algebra}
\externaldocument[brauer-]{brauer}
\externaldocument[homology-]{homology}
\externaldocument[derived-]{derived}
\externaldocument[simplicial-]{simplicial}
\externaldocument[more-algebra-]{more-algebra}
\externaldocument[smoothing-]{smoothing}
\externaldocument[modules-]{modules}
\externaldocument[sites-modules-]{sites-modules}
\externaldocument[injectives-]{injectives}
\externaldocument[cohomology-]{cohomology}
\externaldocument[sites-cohomology-]{sites-cohomology}
\externaldocument[dga-]{dga}
\externaldocument[dpa-]{dpa}
\externaldocument[sdga-]{sdga}
\externaldocument[hypercovering-]{hypercovering}
\externaldocument[schemes-]{schemes}
\externaldocument[constructions-]{constructions}
\externaldocument[properties-]{properties}
\externaldocument[morphisms-]{morphisms}
\externaldocument[coherent-]{coherent}
\externaldocument[divisors-]{divisors}
\externaldocument[limits-]{limits}
\externaldocument[varieties-]{varieties}
\externaldocument[topologies-]{topologies}
\externaldocument[descent-]{descent}
\externaldocument[perfect-]{perfect}
\externaldocument[more-morphisms-]{more-morphisms}
\externaldocument[flat-]{flat}
\externaldocument[groupoids-]{groupoids}
\externaldocument[more-groupoids-]{more-groupoids}
\externaldocument[etale-]{etale}
\externaldocument[chow-]{chow}
\externaldocument[intersection-]{intersection}
\externaldocument[pic-]{pic}
\externaldocument[weil-]{weil}
\externaldocument[adequate-]{adequate}
\externaldocument[dualizing-]{dualizing}
\externaldocument[duality-]{duality}
\externaldocument[discriminant-]{discriminant}
\externaldocument[derham-]{derham}
\externaldocument[local-cohomology-]{local-cohomology}
\externaldocument[algebraization-]{algebraization}
\externaldocument[curves-]{curves}
\externaldocument[resolve-]{resolve}
\externaldocument[models-]{models}
\externaldocument[functors-]{functors}
\externaldocument[equiv-]{equiv}
\externaldocument[pione-]{pione}
\externaldocument[etale-cohomology-]{etale-cohomology}
\externaldocument[proetale-]{proetale}
\externaldocument[relative-cycles-]{relative-cycles}
\externaldocument[more-etale-]{more-etale}
\externaldocument[trace-]{trace}
\externaldocument[crystalline-]{crystalline}
\externaldocument[spaces-]{spaces}
\externaldocument[spaces-properties-]{spaces-properties}
\externaldocument[spaces-morphisms-]{spaces-morphisms}
\externaldocument[decent-spaces-]{decent-spaces}
\externaldocument[spaces-cohomology-]{spaces-cohomology}
\externaldocument[spaces-limits-]{spaces-limits}
\externaldocument[spaces-divisors-]{spaces-divisors}
\externaldocument[spaces-over-fields-]{spaces-over-fields}
\externaldocument[spaces-topologies-]{spaces-topologies}
\externaldocument[spaces-descent-]{spaces-descent}
\externaldocument[spaces-perfect-]{spaces-perfect}
\externaldocument[spaces-more-morphisms-]{spaces-more-morphisms}
\externaldocument[spaces-flat-]{spaces-flat}
\externaldocument[spaces-groupoids-]{spaces-groupoids}
\externaldocument[spaces-more-groupoids-]{spaces-more-groupoids}
\externaldocument[bootstrap-]{bootstrap}
\externaldocument[spaces-pushouts-]{spaces-pushouts}
\externaldocument[spaces-chow-]{spaces-chow}
\externaldocument[groupoids-quotients-]{groupoids-quotients}
\externaldocument[spaces-more-cohomology-]{spaces-more-cohomology}
\externaldocument[spaces-simplicial-]{spaces-simplicial}
\externaldocument[spaces-duality-]{spaces-duality}
\externaldocument[formal-spaces-]{formal-spaces}
\externaldocument[restricted-]{restricted}
\externaldocument[spaces-resolve-]{spaces-resolve}
\externaldocument[formal-defos-]{formal-defos}
\externaldocument[defos-]{defos}
\externaldocument[cotangent-]{cotangent}
\externaldocument[examples-defos-]{examples-defos}
\externaldocument[algebraic-]{algebraic}
\externaldocument[examples-stacks-]{examples-stacks}
\externaldocument[stacks-sheaves-]{stacks-sheaves}
\externaldocument[criteria-]{criteria}
\externaldocument[artin-]{artin}
\externaldocument[quot-]{quot}
\externaldocument[stacks-properties-]{stacks-properties}
\externaldocument[stacks-morphisms-]{stacks-morphisms}
\externaldocument[stacks-limits-]{stacks-limits}
\externaldocument[stacks-cohomology-]{stacks-cohomology}
\externaldocument[stacks-perfect-]{stacks-perfect}
\externaldocument[stacks-introduction-]{stacks-introduction}
\externaldocument[stacks-more-morphisms-]{stacks-more-morphisms}
\externaldocument[stacks-geometry-]{stacks-geometry}
\externaldocument[moduli-]{moduli}
\externaldocument[moduli-curves-]{moduli-curves}
\externaldocument[examples-]{examples}
\externaldocument[exercises-]{exercises}
\externaldocument[guide-]{guide}
\externaldocument[desirables-]{desirables}
\externaldocument[coding-]{coding}
\externaldocument[obsolete-]{obsolete}
\externaldocument[fdl-]{fdl}
\externaldocument[index-]{index}

% Theorem environments.
%
\theoremstyle{plain}
\newtheorem{theorem}[subsection]{Theorem}
\newtheorem{proposition}[subsection]{Proposition}
\newtheorem{lemma}[subsection]{Lemma}

\theoremstyle{definition}
\newtheorem{definition}[subsection]{Definition}
\newtheorem{example}[subsection]{Example}
\newtheorem{exercise}[subsection]{Exercise}
\newtheorem{situation}[subsection]{Situation}

\theoremstyle{remark}
\newtheorem{remark}[subsection]{Remark}
\newtheorem{remarks}[subsection]{Remarks}

\numberwithin{equation}{subsection}

% Macros
%
\def\lim{\mathop{\mathrm{lim}}\nolimits}
\def\colim{\mathop{\mathrm{colim}}\nolimits}
\def\Spec{\mathop{\mathrm{Spec}}}
\def\Hom{\mathop{\mathrm{Hom}}\nolimits}
\def\Ext{\mathop{\mathrm{Ext}}\nolimits}
\def\SheafHom{\mathop{\mathcal{H}\!\mathit{om}}\nolimits}
\def\SheafExt{\mathop{\mathcal{E}\!\mathit{xt}}\nolimits}
\def\Sch{\mathit{Sch}}
\def\Mor{\mathop{\mathrm{Mor}}\nolimits}
\def\Ob{\mathop{\mathrm{Ob}}\nolimits}
\def\Sh{\mathop{\mathit{Sh}}\nolimits}
\def\NL{\mathop{N\!L}\nolimits}
\def\CH{\mathop{\mathrm{CH}}\nolimits}
\def\proetale{{pro\text{-}\acute{e}tale}}
\def\etale{{\acute{e}tale}}
\def\QCoh{\mathit{QCoh}}
\def\Ker{\mathop{\mathrm{Ker}}}
\def\Im{\mathop{\mathrm{Im}}}
\def\Coker{\mathop{\mathrm{Coker}}}
\def\Coim{\mathop{\mathrm{Coim}}}

% Boxtimes
%
\DeclareMathSymbol{\boxtimes}{\mathbin}{AMSa}{"02}

%
% Macros for moduli stacks/spaces
%
\def\QCohstack{\mathcal{QC}\!\mathit{oh}}
\def\Cohstack{\mathcal{C}\!\mathit{oh}}
\def\Spacesstack{\mathcal{S}\!\mathit{paces}}
\def\Quotfunctor{\mathrm{Quot}}
\def\Hilbfunctor{\mathrm{Hilb}}
\def\Curvesstack{\mathcal{C}\!\mathit{urves}}
\def\Polarizedstack{\mathcal{P}\!\mathit{olarized}}
\def\Complexesstack{\mathcal{C}\!\mathit{omplexes}}
% \Pic is the operator that assigns to X its picard group, usage \Pic(X)
% \Picardstack_{X/B} denotes the Picard stack of X over B
% \Picardfunctor_{X/B} denotes the Picard functor of X over B
\def\Pic{\mathop{\mathrm{Pic}}\nolimits}
\def\Picardstack{\mathcal{P}\!\mathit{ic}}
\def\Picardfunctor{\mathrm{Pic}}
\def\Deformationcategory{\mathcal{D}\!\mathit{ef}}
\usepackage{cleveref}
\crefname{originallemma}{original lemma}{original lemmas}
\Crefname{originallemma}{Original lemma}{Original lemmas}
\makeatletter
\newlabel{algebra-lemma-completion-Noetherian}{{05GH}{source}{}{sourcepointer-1}{}}
\newlabel{algebra-lemma-completion-Noetherian@cref}{{[originallemma][0][]05GH}{[1][0][]source}{}{}{}}
\makeatother
\crefname{originallemma}{original lemma}{original lemmas}
\Crefname{originallemma}{Original lemma}{Original lemmas}
\makeatletter
\newlabel{algebra-lemma-hathat-finitely-generated}{{05GG}{source}{}{sourcepointer-2}{}}
\newlabel{algebra-lemma-hathat-finitely-generated@cref}{{[originallemma][0][]05GG}{[1][0][]source}{}{}{}}
\makeatother
\crefname{originallemma}{original lemma}{original lemmas}
\Crefname{originallemma}{Original lemma}{Original lemmas}
\makeatletter
\newlabel{lemma-topologically-nilpotent}{{0AMW}{source}{}{sourcepointer-3}{}}
\newlabel{lemma-topologically-nilpotent@cref}{{[originallemma][0][]0AMW}{[1][0][]source}{}{}{}}
\makeatother
\title{Properties of the completed tensor product of linearly topologized rings}
\author{The Stacks Project Authors}
\date{stacks-project:formal-spaces.tex:local-snapshot-20261012; GNU FDL 1.2 or later}
\begin{document}
\maketitle
\noindent\textit{Editorial note.} All delivered mathematical text is the author's own, copied from the pinned Stacks Project chapter file formal-spaces.tex; no statement or proof step was rewritten, reordered or removed.\par\medskip
\noindent\textit{Editorial note.} Part (5) of the proof contains the author's parenthetical '(details omitted)' attached to the claim that the ideals of topologically nilpotent elements are ideals of definition in the adic Noetherian case; the parenthetical is reproduced verbatim and the verification is recorded as an obligation of the card.\par\medskip
\noindent\textit{Editorial note.} The proof displays the definition of the completed tensor product as a limit, the containment between powers of the kernels and the identification of a quotient with a tensor product of quotients; all lie inside the counted span and are reproduced verbatim.\par\medskip
\noindent\textit{Editorial note.} The counted spans carry 3 navigation stubs; no bibliography citation occurs in the spans.\par\medskip
\noindent\textit{Editorial note.} The preamble span points at the chapter's own preamble.tex shipped with the frozen source, used verbatim as staged.\par\medskip
\noindent\textit{Allowed standard prerequisites (editorial).}\begin{itemize}
\item Weakly admissible, admissible, pre-adic and adic rings with their ideals of definition. Same chapter, the definitions of these notions, used throughout at proof lines 1223-1245 and 1263-1282. Supply the properties being transferred to the completed tensor product.
\item The closure of the ideal generated by two systems of topologically nilpotent elements consists of topologically nilpotent elements. Same chapter, Lemma lemma-topologically-nilpotent, invoked at proof lines 1226-1228. Gives part (1).
\item Completion with respect to a finitely generated ideal, and Noetherianness of a completion. Algebra, Lemma algebra-lemma-hathat-finitely-generated at proof lines 1258-1260, and Algebra, Lemma algebra-lemma-completion-Noetherian at proof lines 1281-1282. Give parts (4) and (5) respectively.
\end{itemize}\medskip
\ifdefined\hypertarget\hypertarget{sourcepointer-1}{}\fi
\noindent\textit{Source reference.} Original lemma 05GH: sources/stacks-chow/tags.txt, source lines 7061--7061. Pointer outside the counted span, invoked once for context; the load-bearing argument of the counted statement is the author's proof reproduced in full in the proof span.\par\medskip
\ifdefined\hypertarget\hypertarget{sourcepointer-2}{}\fi
\noindent\textit{Source reference.} Original lemma 05GG: sources/stacks-chow/tags.txt, source lines 7060--7060. Pointer outside the counted span, invoked once for context; the load-bearing argument of the counted statement is the author's proof reproduced in full in the proof span.\par\medskip
\ifdefined\hypertarget\hypertarget{sourcepointer-3}{}\fi
\noindent\textit{Source reference.} Original lemma 0AMW: sources/stacks-formal-spaces/pinned-original-source.tex, source lines 1134--1134. Pointer outside the counted span, invoked once for context; the load-bearing argument of the counted statement is the author's proof reproduced in full in the proof span.\par\medskip
\begin{lemma}
\label{lemma-completed-tensor-product}
Let $B \to A$ and $B \to C$ be continuous homomorphisms of
linearly topologized rings.
\begin{enumerate}
\item If $A$ and $C$ are weakly pre-admissible, then
$A \widehat{\otimes}_B C$ is weakly admissible.
\item If $A$ and $C$ are pre-admissible, then
$A \widehat{\otimes}_B C$ is admissible.
\item If $A$ and $C$ have a countable fundamental system of open
ideals, then $A \widehat{\otimes}_B C$ has a countable fundamental
system of open ideals.
\item If $A$ and $C$ are pre-adic and have finitely generated ideals
of definition, then $A \widehat{\otimes}_B C$ is adic and has
a finitely generated ideal of definition.
\item If $A$ and $C$ are pre-adic Noetherian rings and
$B/\mathfrak b \to A/\mathfrak a$ is of finite type
where $\mathfrak a \subset A$ and $\mathfrak b \subset B$
are the ideals of topologically nilpotent elements, then
$A \widehat{\otimes}_B C$ is adic Noetherian.
\end{enumerate}
\end{lemma}

\begin{proof}
Let $I_\lambda \subset A$, $\lambda \in \Lambda$ and
$J_\mu \subset C$, $\mu \in M$
be fundamental systems of open ideals, then by definition
$$
A \widehat{\otimes}_B C =
\lim_{\lambda, \mu} A/I_\lambda \otimes_B C/J_\mu
$$
with the limit topology. Thus a fundamental system of open ideals
is given by the kernels $K_{\lambda, \mu}$ of the maps
$A \widehat{\otimes}_B C \to A/I_\lambda \otimes_B C/J_\mu$.
Note that $K_{\lambda, \mu}$ is the closure of the ideal
$I_\lambda(A \widehat{\otimes}_B C) + J_\mu(A \widehat{\otimes}_B C)$.
Finally, we have a ring homomorphism
$\tau : A \otimes_B C \to A \widehat{\otimes}_B C$ with dense image.

\medskip\noindent
Proof of (1). If $I_\lambda$ and $J_\mu$ consist of topologically
nilpotent elements, then so does $K_{\lambda, \mu}$ by
Lemma \ref{lemma-topologically-nilpotent}. Hence 
$A \widehat{\otimes}_B C$ is weakly admissible by definition.

\medskip\noindent
Proof of (2). Assume for some $\lambda_0$ and $\mu_0$ the ideals
$I = I_{\lambda_0} \subset A$ and $J_{\mu_0} \subset C$ are ideals of
definition. Thus for every $\lambda$ there exists an $n$ such that
$I^n \subset I_\lambda$. For every $\mu$ there exists an $m$ such that
$J^m \subset J_\mu$. Then
$$
\left(I(A \widehat{\otimes}_B C) + J(A \widehat{\otimes}_B C)\right)^{n + m}
\subset
I_\lambda(A \widehat{\otimes}_B C) + J_\mu(A \widehat{\otimes}_B C)
$$
It follows that the open ideal $K = K_{\lambda_0, \mu_0}$
satisfies $K^{n + m} \subset K_{\lambda, \mu}$. Hence $K$
is an ideal of definition of $A \widehat{\otimes}_B C$
and $A \widehat{\otimes}_B C$ is admissible by definition.

\medskip\noindent
Proof of (3). If $\Lambda$ and $M$ are countable, so is
$\Lambda \times M$.

\medskip\noindent
Proof of (4). Assume $\Lambda = \mathbf{N}$ and $M = \mathbf{N}$
and we have finitely generated ideals $I \subset A$ and $J \subset C$
such that $I_n = I^n$ and $J_n = J^n$. Then
$$
I(A \widehat{\otimes}_B C) + J(A \widehat{\otimes}_B C)
$$
is a finitely generated ideal and it is easily seen that
$A \widehat{\otimes}_B C$ is the completion of
$A \otimes_B C$ with respect to this ideal. Hence (4)
follows from Algebra, Lemma \ref{algebra-lemma-hathat-finitely-generated}.

\medskip\noindent
Proof of (5).
Let $\mathfrak c \subset C$ be the ideal of topologically nilpotent elements.
Since $A$ and $C$ are adic Noetherian, we see that
$\mathfrak a$ and $\mathfrak c$ are ideals of definition (details omitted).
From part (4) we already know that
$A \widehat{\otimes}_B C$ is adic and that
$\mathfrak a(A \widehat{\otimes}_B C) + \mathfrak c(A \widehat{\otimes}_B C)$
is a finitely generated ideal of definition. Since
$$
A \widehat{\otimes}_B C /
\left(\mathfrak a(A \widehat{\otimes}_B C) +
\mathfrak c(A \widehat{\otimes}_B C)\right)
=
A/\mathfrak a \otimes_{B/\mathfrak b} C/\mathfrak c
$$
is Noetherian as a finite type algebra over the Noetherian ring
$C/\mathfrak c$ we conclude by
Algebra, Lemma \ref{algebra-lemma-completion-Noetherian}.
\end{proof}

\end{document}
